% /root/monodromy, 2026-11-19. Both local proofs independently audited by root and RD; no global extension asserted.
\subsection{The corresponding local criterion at an ordinary triple point}
\label{subsec:ss-triple-point}
\status{Proved under the stated hypotheses. The parent and reduction/descent agent independently audited both complete local proofs. No global incidence or full-monodromy assertion is made.}

\begin{lemma}[Smoothness at a mixed ordinary triple point]
\label{lem:ss-mixed-triple-point}
Let $A',A$ be superspecial abelian surfaces in characteristic two,
and let $\phi:A'\to A$ be an odd-degree isogeny. Let $C'\subset A'$
be smooth and suppose $\phi|_{C'}$ is birational onto its image $Z$.
Assume $Z$ has ordinary nodes and exactly one ordinary triple point,
with preimages $P_1,P_2,P_3$. Choose tangent parameters on these
three branches such that their tangent vectors, identified by
$\phi$, satisfy $v_3=v_1+v_2$.
Suppose
\[
 \ker C_{C'}=H^0(A',\Omega^1_{A'})|_{C'}\oplus k\eta,
 \qquad \eta(P_1)+\eta(P_2)+\eta(P_3)\ne0,
\]
where values use the chosen tangent vectors. Suppose
$\mathcal N=\mathcal O_A(Z)$ is globally generated and its
polarization is separable. Then the relative coherent Cartier map
has corank one at $Z$, and its determinant has a nonzero first
derivative there. The relative Hasse divisor is therefore smooth
at $Z$.
\end{lemma}
\begin{proof}
The three-point version of Lemma~\ref{lem:ss-two-point-Cartier}
is obtained by replacing $P+Q$ with $P_1+P_2+P_3$ in its proof.
For completeness, the zero-residue subspace of
$H^0(\omega_{C'}(2P_1+2P_2+2P_3))$ has dimension $g(C')+3$,
and its quotient by holomorphic forms is the space $k^3$ of
double-pole coefficients. Its Cartier map to the holomorphic
Cartier cokernel is dual, up to the scalar Frobenius twists, to
\[
     \ker C_{C'}\longrightarrow k^3,
     \quad\xi\longmapsto(\xi(P_1),\xi(P_2),\xi(P_3)).
\]
Indeed the constant-difference kernel in the twisted Frobenius
calculation is now $k^3/k$. Frobenius is bijective on it; after
its unique adjustment, subtracting a common constant kills all
three zeroth-order terms. The three first-jet terms of the Cech
primitives are $\xi(P_i)$, and the local residue pairing gives
the displayed map. This repeats the complete two-point calculation
without any change to its duality hypotheses.

We spell out the ordinary-triple-point matching. Choose completed
coordinates so that the first two branches are $y=0,x=0$, and
the third is $y=\varphi(x)$, where
$\varphi(x)=x+O(x^2)$. A local equation is
\[
                         f=xy(y-\varphi(x)).
\]
The three adjunction generators have leading terms
$dx/x^2$, $dy/y^2$, and $dx/x^2$ respectively.
Thus all three double-pole coefficients of a dualizing form are
equal. Their residues sum to zero. These are also sufficient:
multiplication of the adjunction generator by the constant and
the two linear terms of a numerator produces the full
three-dimensional space of allowed polar parts. Equivalently,
the conductor is the square of the maximal ideal on each
normalized branch; its polar conditions are equality of the
three leading coefficients and the single residue relation.

Cartier-killed forms have zero residues. The only possible polar
vector at the triple point is therefore a multiple of $(1,1,1)$.
Invariant ambient forms have value sum zero since $v_1+v_2+v_3=0$.
The hypothesis on $\eta$ and the three-point calculation exclude
a nonzero such polar vector. Ordinary nodes likewise contribute
no pole to a Cartier-killed form. Consequently
\[
               \ker C_Z=\ker C_{C'},\qquad\dim\ker C_Z=3.
\]
The adjunction boundary on the two invariant forms is the
differential of $\lambda_{\mathcal N}$, after contraction with
an invariant volume form. Its explicit Cech calculation is in
Lemma~\ref{lem:ss-mixed-self-tacnode}. Separability makes it an
isomorphism. Thus a unique ambient correction of $\eta$ lies in
the relative Cartier kernel; denote it by $\eta_0$. Its three-value
sum is still nonzero, and that relative kernel has dimension one.

Lift $\eta_0$ to $u\in W=H^0(A,\mathcal N)$ modulo $ks$, where
$s$ defines $Z$. Holomorphicity on all branches forces the
constant and both linear terms of $u$ to vanish. Write its
quadratic part as
\[
                          u=a x^2+bxy+c y^2.
\]
After normalizing the invariant volume form at the singular point,
the branch values are $a,c,a+b+c$. Their sum is $b$, hence $b\ne0$.
The equation $s$ has order three. Therefore the Cartier trace of
$sv$ has zero value at the singular point for every regular $v$,
since that value extracts the $xy$ coefficient.

In root-removed coordinates the relative map is
\[
                  M(s):W/ks\longrightarrow W^{(2)}/ks^{(2)}.
\]
Frobenius-twisted evaluation at the singular point spans its
one-dimensional left kernel. Global generation supplies a
section $\dot s$ whose value there is one. The $xy$ coefficient
of $\dot s\,u$ is $b\ne0$, so
\[
                \operatorname{ev}^{(2)}(M(\dot s)u)\ne0.
\]
As in the self-tacnode lemma, derivatives of the input lift are
killed by the left kernel, and the target line $ks^{(2)}$ has no
first variation in characteristic two. The adjugate derivative
formula proves that the determinant has a nonzero derivative.
All quantities are intrinsic on the projective coefficient space:
the radial derivative vanishes at the determinant-zero point.
\end{proof}

\begin{remark}[A remaining global, rather than local, question]
The hypothesis on the three values has not been proved on every
triple-identification divisor in a norm family. In particular,
triples lying on a single elliptic fibre can have a different
incidence dimension from three general points. This lemma does
not discard those components and makes no assertion that the
needed value sum is generically nonzero on them.
\end{remark}

\begin{lemma}[A tangent triple joining two collision types]
\label{lem:ss-tangent-triple-bridge}
In the same superspecial-isogeny setting, suppose the normalization
$C'$ is sharp in the sense that
\[
             \ker C_{C'}=H^0(A',\Omega^1_{A'})|_{C'}.
\]
Suppose its image $Z$ has ordinary nodes and one further singular
point with three smooth branches: the first two have simple contact,
and the third is transverse to their common tangent. Assume
$\mathcal N=\mathcal O_A(Z)$ is very ample and its polarization is
separable. Then the relative Hasse divisor is smooth at $Z$.
\end{lemma}
\begin{proof}
Let $P,Q,R$ be the preimages of the exceptional singularity, with
$P,Q$ the tangent pair. Choose completed coordinates so that the
branches are $y=0$, $y=\varphi(x)$, and $x=0$, with
$\varphi(x)=c x^2+O(x^3)$ and $c\ne0$. A local equation is
\[
                         f=xy(y-\varphi(x)).
\]
The adjunction generators have pole orders three on the first
two branches and order two on the third. A Cartier-killed
meromorphic differential has no term $x^{-3}dx$, because Cartier
maps it to a nonzero multiple of $x^{-2}dx$; it also has zero
residue. Thus the only possible nonzero polar terms are double
poles at $P,Q,R$.

The dualizing conditions for those zero-residue terms reduce to
\[
                     (a_P,a_Q,a_R)=(a,a,0).
\]
One can read this directly from adjunction: multiplication of
$dx/(x\varphi(x))$ by a numerator with leading term $x$ gives
matching double-pole coefficients at $P,Q$ and none on the third
branch. A quadratic correction to that numerator removes the
matching residues. Conversely multiplication of a dualizing form
by either local coordinate must have residues summing to zero.
The coordinate $x$ gives $a_P+a_Q=0$, and the coordinate $y$
gives $a_R=0$. These are all conditions on zero-residue double
terms: their one-dimensional solution space is realized by the
numerator just described. Holomorphic terms are unrestricted.

By the three-point Cartier boundary calculation, a killed form
with polar vector $(1,1,0)$ exists precisely when
$\xi(P)+\xi(Q)=0$ for every $\xi\in\ker C_{C'}$.
This holds because the kernel consists of invariant ambient
forms and the two tangents agree under translation. The killed
meromorphic forms with zero polar vector are exactly the two
ambient holomorphic forms. Hence
\[
                        \dim\ker C_Z=3.
\]
The separable-polarization boundary is an isomorphism on those
two ambient forms. After their unique subtraction the relative
Cartier kernel is therefore one-dimensional, and its generator
still has nonzero double-pole coefficients on the tangent pair.

Let $u\in H^0(\mathcal N)$ be a numerator lifting that generator.
The preceding local matching shows that $u$ has nonzero linear
$x$ term and zero constant and linear $y$ terms. The defining
section $s$ has order three, so Frobenius-twisted evaluation at
the singular point annihilates the entire Cartier image and
spans its relative left kernel. Very ampleness supplies a global
$\dot s$ with zero value and first jet $y$. Then the $xy$
coefficient of $\dot s\,u$ is nonzero. The same quotient
adjugate derivative calculation as above proves that the relative
Hasse determinant has nonzero derivative at $Z$.
\end{proof}

\begin{remark}[What the tangent-triple bridge would supply]
A norm-family parameter with this singularity belongs both to a
pair-tangency locus and to a triple-identification locus. If it is
a smooth point of the relevant target Hasse divisor, any target
Hasse components containing the images of these loci coincide.
The lemma proves that local smoothness under its stated hypotheses.
It does not prove the existence or dominance of such parameters,
or the irreducibility of the source collision loci. Those remain
separate global incidence questions.
\end{remark}
